
\section{怎样写比较复杂的记叙文}

《中学语文教学大纲》对高中一年级学生提出写作训练的要求是：能写“比较复杂的记叙文”，做到“观点鲜明，内容充实，结构完整，条理清楚，语句流畅。”其实，记叙文的基本写法已如上述。这个“比较复杂”只是跟“比较简单”相对而言的。究竟“比较复杂”在哪些地方呢？大致说来，包括以下各个方面：记叙的往往不止是一个人、一件事、一个场景；所涉及的人物之间的关系比较错综复杂，即使是只写一个人，也得涉及这个人生活的各个方面，需要表现出这个人思想性格的曲折变化；所记的几件事之间的关系也比较错综复杂，即使是只写一件事，这件事延续的时间一定比较长，发展变化一定比较曲折；所描写的场景不断转换，即使是只写一景，也得从不同的角度或不同的条件来写，使一景变成多景。由于头绪比较纷繁，涉及的人或事比较多，历时比较长，生活面比较广，我们就把这一类记叙文叫做“比较复杂的记叙文”。

怎样来写比较复杂的记叙文呢？应注意哪些问题呢？

\textbf{1. 用几件事写人，要注意紧扣中心，选材精当。}
写比较复杂的记叙文，不能想到一个人、一件事就够了，而要多方面展开联想和想象，把与题目有关的人和事尽量搜拢来；然后根据确定的中心意思，选取最能突出主题的部分，有层次有重点地加以叙述和描写，进行精心的剪裁和组织。例如课文《谁是最可爱的人？》，为了深刻地表达志愿军战士的确是最可敬爱的人这个中心思想，作者魏巍从他搜集到的大量的素材中只选取了三个事例——松骨峰战斗、战士马玉祥冒火抢救朝鲜儿童，一个战士和作者一段满怀深情的谈话。这三个事例分别写出志愿军战士对敌人的刻骨仇恨，对朝鲜人民的无比热爱，对祖国对人民的无限忠诚，从而显示出志愿军战士崇高的革命英雄主义、国际主义和爱国主义精神，三者集中起来，就最热烈地歌颂了志愿军战士是我们最可爱的人。可以看出，作者从大量素材中提炼出来的主题是鲜明突出的，他所选取的材料是紧扣中心思想的，而且他所选用的材料都能表现出众多素材中具有共性的思想，各个方面的材料之间有严密的内在联系，能共同表达一个深刻的主题。有的同学在写这一类比较复杂的记叙文时，往往会犯这样的毛病：一是还没有提炼好鲜明的中心思想，就匆匆下笔，因此把想到的有关素材一古脑儿写出来，使文章杂乱无章，不知所云；或者所写的材料各不相关没有内在联系，不能表现一个共同的意思。二是在有了明确的中心以后，行文时却不能紧扣中心，往往不知不觉地往外岔；或者对自己想到的某件事感到生动有趣，不管能不能表现中心思想也硬往上凑，不忍割爱。这些都是应尽力避免的。

写人所用的几件事，可以用来表现某人的一种品质，也可以用来表现某人的几种品质。例如课文《人民的勤务员》，作者选取了“给大嫂买车票”、“帮助大娘找儿子”、“冒雨送大嫂母子回家”、“火车上做好事”、“打扫候车室”五个事例，从不同的方面表现了雷锋同志全心全意为人民服务的精神。这种写法的好处是内容集中，能表现出一定的深度。而象《谁是最可爱的人》，则通过三件事，表现了志愿军战士多方面的崇高品质和伟大精神。这种写法的好处，是能比较全面地表现出一个人或一类人的思想品质和精神面貌。我们可以根据具体情况学习运用。

\textbf{2. 写几个人，要注意突出写中心人物，并写好这个中心人物与其他人物的关系。}

有些以写人为主的记叙文题目，如《我家的四邻》、《他们都有美好的心灵》、《接力赛中显英豪》、《我们可爱的班集体》等，要求我们写的不只是一个人，而是好多人。在短短的一篇作文里，当然不可能把所有涉及的人都详细写出来，只能选取一个最突出、最有代表性的人作重点描写，通过这个中心人物与其他人物之间的关系，就能反映出一定的社会现实，表达一个鲜明的中心思想。以课文《坚强的战士》为例，描写的中心人物是共产党员林红（就《青春之歌》这部长篇小说而言，它的中心人物是林道静），作者杨沫着重刻画了林红如何继承丈夫李伟的遗志，教育和鼓舞狱中的难友、革命青年林道静和俞淑秀，直至英勇就义，从而引导林道静和俞淑秀走上了革命的道路。这样写，不仅歌颂了共产党员林红坚强不屈的性格和献身革命的崇高品质，而且深刻地表达了“共产主义是不可战胜的”这一光辉的思想。如果作文题是叫我们写集体活动或场面，涉及的人一般比较多，但是记叙的重点是事，而不是人，那就应该抓住事件、活动、场面的有特色的东西来写，其中的人物不必一一点名，逐个描写，只要选取一两个有代表性的突出人物写一下，能反映出事件、活动、场面的面貌就行了。

\textbf{3. 要综合运用各种有效的记叙方法，安排好材料。}

既然比较复杂的记叙文涉及的人或事比较多，头绪比较纷繁，我们就需用相应的比较复杂的形式来表达，才能做到观点鲜明，内容充实，结构完整，条理清楚。

首先，我们要运用点面结合的方法，使重点材料与一般性材料、重点的局部情况与一般的总体情况结合起来。比如写一件事，我们既着力写那些最能体现事件意义的片断或侧面，又简要交代事件的发展过程；写一个人，既具体记叙最能表现他思想性格的一件或几件事，又对他的概况作些必要的介绍；写一个场面，既有对个别人或事的特写，又有对全景的鸟瞰。这样把点面材料安排好了，才能以一般映衬个别，相辅相成地表现出生活本身固有的丰富多彩，使文章记叙完整，重点突出，详略得当，语言简练。如果文章只有“面”而无“点”，容易流于笼统平淡，文字拖沓；相反，只有“点”而无“面”，则会失之偏狭单薄，难窥全貌。

运用点面结合的方法处理材料，应该注意，

（1）哪些材料属于“点”，哪些材料属于“面”，取决于它们对表现中心思想的作用的大小。

（2）在行文中，“面”上的材料要写得较为概括简略，“点”上的材料要写得比较细致具体。

（3）文章由“点”到“面”，或由“面”到“点”，都要发展得自然无痕，不要互相游离脱节。

其次，我们要努力为文章找一条合理的线索。所谓线索，是指在记叙文中能把全部材料串成一个有机整体的脉络，它起着组织情节、结构全文的作用。有了它，我们就可以替各种材料找到适当的位置，决定哪些该在前，哪些该在后，并把各种材料串联起来。

可以作为线索的东西是多种多样的，时间、地点、物件以至人物等等都可以。以什么作为一篇文章的线索，需要我们慎重地加以选择。选择的原则是：作为线索的事物，必须能体现材料之间的内在联系，有利于突出中心思想。例如课文奥斯特洛夫斯基的《我的一天》，用时间推移作线索，有利于清楚地反映出“我”在清晨、上午、下午都做了些什么，从而表现“我”作为一个革命战士的坚强和幸福；鲁迅的《从百草园到三味书屋》，用空间转换作为线索，能把有关百草园的趣事和三味书屋压抑沉闷的景况作鲜明的对比；谢璞的《珍珠赋》，用珍珠作为文章的线索，在于“珍珠”既可以用来指实物，又可以用来作比方，能够把“虚”和“实”结合起来，起到统帅全部材料、点明中心思想的作用；曹靖华的《小米的回忆》，则用小米这种粮食为线索，不仅把童年、一九三四年和抗日战争时期难忘的人和事串联了起来，而且突出了“延安精神”是我们战胜一切困难、夺取革命胜利的力量源泉这一深刻的中心思想。我们在练习写作比较复杂的记叙文时，应该学习这些范文选择线索的方法，使文章能做到一线贯串，中心突出。

再次，还应交叉运用多种叙述顺序，把最能表现中心思想的若干材料巧妙地、有机地组织起来，以反映丰富多彩的社会生活。古人说：“文如看山不喜平。”平铺直叙，记流水账，往往使文章平淡寡味，不耐咀嚼。因此，我们除了可以采取顺序结合插叙和顺叙结合倒叙这两种比较简单易的形式之外，还可以采取先倒叙、后顺叙、中间结合运用插叙（或补叙）这样比较复杂的形式，也可以在顺叙中多次使用插叙（或补叙）。当然，究竟采取哪种叙述方式为好，要根据表达中心思想的需要来决定，并没有什么固定的模式。在交叉运用多种叙述方式时，要注意：一定要使各部分“接榫”的地方衔接紧密，并使各部分的起止界线一清二楚；一定要使故事的主线突出、清晰、贯串始终，不能因插叙过多而造成喧宾夺主。如果过分追求形式上的曲折起伏而故弄玄虚，搞得文章头绪不清，缺乏条理，甚至中心不明，那就舍本逐末了。

\textbf{4. 要综合运用多种描写手法，力求文章形象生动。}

刻画人物形象，总离不开对人物的描写。描写人物可以直接地正面描写，也可以间接地侧面描写。正面描写是对人物的外貌、语言、神态、动作、心理活动等方面作直接的描绘。在写比较复杂的记叙文时，我们还可以对人物的这些方面的细小环节作细致、具体的描写，这叫做细节描写。好的细节描写不仅能增强文章的生动性和真实感，而且对刻画人物的性格和表现中心思想有一定的作用。至于侧面描写，则是对人物不作直接刻画，而是通过对与他有联系的其它人物或景物的描写，间接烘托出作者所要着力描写的人物的性格和思想。这种描写手法，可使文章含蓄，耐人寻味，还能烘托气氛，使文章活泼、生动，使人物的精神面貌更加突出。

侧面描写常采用以下几种方法：
（1）通过描述别人的对话来进行。如鲁迅小说《药》，并未直接描写革命者夏瑜，只是通过康大叔和一群愚昧守旧的人的谈话，突出烈士的革命事迹和烈士的言行不被一般群众理解的情景，从而既使夏瑜的英姿跃然纸上，又揭示了资产阶级民主革命严重脱离群众的悲剧根源。

（2）通过别人的表情、神态、动作的描写来衬托。如赵树理的《小二黑结婚》，写“小芹去洗衣服，马上青年们也都会去洗，小芹上树采野菜，马上青年们也都去采。”从一群小伙子的表现上，侧面烘托了小芹的漂亮出众，着墨不多，效果却很好。

（3）通过景物描写来衬托。如老舍的《在烈日和暴雨下》（选自《骆驼祥子》）一文，描写“日”之“烈”和“雨”之“暴”，正是为了衬托祥子拉车生活的痛苦，表现旧社会人力车夫所过的牛马不如的生活。

要综合运用叙述、描写、议论、抒情、说明等各种表达方式，为刻画人物、深化主题服务。

比较简单的记叙文，虽然有时也要说明事理，发表议论，抒发感情，但一般都比较短，紧接着事件的叙述，不单独成段。比较复杂的记叙文可以有成段的议论、抒情和说明的段落。如课文《谁是最可爱的人》、《包身工》、《风景谈》、《花城》等名作中，都有脍炙人口的议论、抒情或说明的文字。

应该注意：记叙文是“以事感人”的，行文中总要以叙述和描写为主。议论和抒情，说明都要紧紧围绕所写的人和事进行。记叙文中的议论，只是对所记叙的事物发表精辟的见解，点明所写的人或事的深刻意义，常起画龙点睛的作用，因此不必象议论文那样要有完整的论证过程，文字要简明扼要。记叙文中的抒情，则是在记叙的过程中，作者“情动而辞发”，鲜明地表露作者不可压抑的爱情，突出自己的观点。这种感情既可以直接抒发，也可以借景抒发，都要做到健康、真实，“以情动人”，并且紧扣所记叙的事物，作简要、强烈的表述。记叙文中的说明，则多用来介绍文中已经出现或将要出现的事物，使读者对这些不十分熟悉或不十分明白的事物有足够的了解，因此文字更不能繁琐冗长，点到就行。在记叙文写作中，一般是在记叙之后作精警的议论、炽热的抒情或简要的说明，它们之间要结合紧密，水乳交融，富有严密的逻辑联系。例如茅盾的《风景谈》，在主要运用叙述和描写方式写出了六个动人画面之外，还十分贴切、十分简洁地进行了议论、抒情和说明，使全文立意深刻，爱憎鲜明，很有感染力，成为不可多得的佳作。我们应该很好地体会、学习。
\newpage